% Copyright (c) 2026 Twin Vector Labs LLC.
% All rights reserved.
\documentclass[11pt]{article}

\usepackage[margin=1in]{geometry}
\usepackage{amsmath,amssymb,amsthm}
\usepackage{mathtools}
\usepackage{enumitem}
\usepackage{booktabs}
\usepackage{array}
\usepackage[colorlinks=true,linkcolor=blue,urlcolor=blue]{hyperref}

\newtheorem{remark}{Remark}

\newcommand{\C}{\mathbb{C}}
\newcommand{\R}{\mathbb{R}}
\newcommand{\Z}{\mathbb{Z}}
\newcommand{\rep}[1]{\mathbf{#1}}
\newcommand{\norm}[1]{\left\lVert #1 \right\rVert}
\newcommand{\abs}[1]{\left\lvert #1 \right\rvert}
\newcommand{\sigmacol}{\sigma}
\newcommand{\Pin}{P_{\mathrm{in}}}
\newcommand{\Pout}{P_{\mathrm{out}}}
\newcommand{\dd}{\mathrm{d}}

\title{\bfseries The Flavor and Charge Sectors --- Preliminary Results:\\
what the implementations and tests established}
\author{Tessera --- cobordism subsystem (epic \#410: \#412, \#413, \#415, \#417;
\#411/\#416/\#418 in flight)}
\date{}

\begin{document}
\maketitle

\begin{abstract}
This is a status report on the \emph{Flavor and Charge Sectors} exploration track
(epic \#410), reporting the results actually \emph{measured} as its tickets were built
and tested. It is the companion to the plan note \emph{The Flavor and Charge Sectors}
(\#420), which laid out the program; here we record what landed. Four results are merged
and verified. (1) \textbf{Orientation is now geometric} (\#412): the emergent
charge read-out was made relabeling-invariant by applying the surface's induced
orientation symmetrically with the inputs, turning the color charge into the geometry's
Stokes invariant ($\sigmacol_R = 0$ for a neutral input, exact; $3$ for a colored one).
(2) \textbf{The proton singlet is now exact} (\#413): replacing the asymmetric prism
extrusion with a \emph{symmetric apex interior} drives the intertwining residual from
$4.26\times10^{-2}$ to $4.5\times10^{-14}$ and the color-singlet overlap from $0.999$ to
$1.000000$ --- the transport is exactly equivariant --- and it is now the default
construction. A side finding: the spontaneous ``photon'' (a worldline relaxing through
null) does \emph{not} survive the symmetrization; it was triangulation-assisted. (3) The
\textbf{Dirac--K\"ahler operator} (\#415) is assembled and verified: $(\dd+\delta)^2=L$,
the Clifford relations exactly, the conserved current $j^0$, and the four-fold taste
multiplicity. (4) The \textbf{electromagnetic field strength splits} into electric and
magnetic parts by causal type (\#417). Three further tickets (the Gauss-law charge
\#411, the bipartite revisit \#416, the $S^3$ spike \#418) are in flight; the flavor
synthesis (\#414) is pending. All numbers are measured through the production C++ and
its Python test suite; nothing here is projected.
\end{abstract}

\tableofcontents

\section{The program, and what this reports}
The track finds the proton's electroweak content (charge, flavor, spin/current) by
splitting the behavior of the \emph{existing} register --- spatio-temporally and across
form-degree --- and by removing geometric arbitrariness, rather than bolting on parallel
registers. The plan is documented separately (\#420). This note reports the
\emph{measured} outcomes of the first wave of implementations. Each result below was
landed as a pull request with its own tests, merged to \texttt{main}, and reproduced
through the production \texttt{TransportCobordism} / \texttt{EigenstateSynthesis} /
\texttt{DiracKahler} path; the figures are from those tests.

\section{Geometric foundations}

\subsection{Orientation is geometric, not a vertex-label sort (\#412)}
\label{sec:orient}
The color charge of a window is a sum of per-hole periods. The inputs were already read
in the surface's induced orientation (the \texttt{endSignCovector}, a globally coherent
$\pm1$ propagated across shared facets), but the \emph{emergent result} was summed in
each hole's own sorted-vertex orientation --- a function of the arbitrary vertex
numbering. A relabeling could flip an individual hole's term with no conjugating partner,
scrambling the sum. The fix applies the result block's induced signs symmetrically with
the inputs. The effect is decisive: the signed carried charge becomes the geometry's
Stokes invariant.

\begin{center}
\small
\begin{tabular}{@{}lcc@{}}
\toprule
input & signed $\abs{\sigmacol_R}$ (induced orientation) & raw $\abs{\sigmacol_R}$ (sorted) \\
\midrule
neutral ($q\bar q$ triple, $\Sigma=0$) & $0.0000$ (exact) & $0.2258$\\
colored ($R,G,B$, $\Sigma=1$ each)     & $3.0000$ (exact) & $1.0219$\\
\bottomrule
\end{tabular}
\end{center}

\noindent The signed sum is exactly the Stokes charge $\sigmacol_R=-(\sigmacol_A+
\sigmacol_B+\sigmacol_C)$; the raw sum is neither, and is label-dependent. The fix is
verified by a relabeling-invariance test (the track's $G_6$ invariant): under seeded
permutations of the base-surface vertex ids the signed $\sigmacol_R$, singlet overlap,
and intertwining residual are invariant up to a single global sign, while the raw sum
drifts. The companion clarification --- that the read-out's two notions of
``orientation'' are distinct --- led to renaming the Causal-Dynamical-Triangulations
$(t_i,t_f)$ class to \texttt{TemporalOrientation}, reserving the \emph{spatial} induced
orientation for \texttt{ChainComplex::endSignCovector}.

\subsection{The symmetric apex interior: the proton singlet is now exact (\#413)}
\label{sec:apex}
The trivalent junction $W_{ABC}$ was built by extruding the holed base surface
$\times I$ and tetrahedralizing each triangular prism by the Freudenthal ``staircase'' ---
whose diagonal is chosen by a vertex-label sort. That choice is not a symmetry of the
geometry, and it --- not the metric --- was the documented source of the intertwining
residual $\norm{M\Pin-\Pout M}/\norm{M}\approx 4.26\times10^{-2}$ that held the color
singlet at $0.999$.

The replacement is a \emph{symmetric apex stacking}, with \emph{no} vertex-label sort
anywhere. Each base triangle cones up to a single \emph{face-apex} vertex (a $(3,1)$
tetrahedron) and down from the apex to the top copy (a $(1,3)$ tetrahedron); the gap
between two adjacent apex-columns over a shared base edge is an octahedron split along the
\emph{canonical dual edge} $f_1 f_2$ (the two face-centres) into four tetrahedra. Because
the dual edge is fixed by the two incident faces rather than by ids, the subdivision is
equivariant under the icosahedral symmetry. The transport then intertwines the color
$\Z_3$ \emph{exactly}:

\begin{center}
\small
\begin{tabular}{@{}lcc@{}}
\toprule
 & prism (Freudenthal) & \textbf{symmetric apex (now default)}\\
\midrule
intertwining residual $\norm{M\Pin-\Pout M}/\norm{M}$ & $4.26\times10^{-2}$ & $\mathbf{4.5\times10^{-14}}$\\
color-singlet overlap & $0.999$ & $\mathbf{1.000000}$\\
valid manifold (\texttt{dualComplexValid}) / $b_1$ & yes / $11$ & yes / $11$\\
\bottomrule
\end{tabular}
\end{center}

\noindent The construction is exposed as \texttt{Spacetime::symmetricStackCells} and is
the default interior on \texttt{TripartiteRegisterTopology}
(\texttt{set\_symmetric\_interior(false)} selects the legacy prism). The twenty existing
junction tests pass unchanged on the symmetric default --- no re-baselining was needed,
the residual simply moved inside their tolerances --- so the color-confinement skeleton
(charge conservation exact at $\sim10^{-14}$, $\sigmacol\to0$, the $\Z_3$ spectrum
$\{1,\omega,\omega^2\}$) is preserved while the singlet becomes exact.

\begin{remark}[The spontaneous photon was triangulation-assisted]
\label{rem:photon}
On the prism, making the cross-layer worldlines timelike and relaxing drives one
worldline through null ($\ell^2\to0$) --- the predicted ``photon.'' On the symmetric
interior the same seed does \emph{not} produce it: through $120$ relaxation iterations all
$450$ worldline edges stay uniformly timelike at $\ell^2=-0.3$, with \emph{zero} near-null
edges. The symmetric relaxation preserves the worldline symmetry; the prism's photon was
triggered by the triangulation asymmetry singling out one edge. The physical reading is
that a photon should emerge from a genuine symmetry-breaking \emph{source}, not from the
bare uniform seed --- the symmetric interior is the more honest setting, and a
source-driven photon is the natural follow-up. (The Lorentzian seed itself works on both
interiors: a structural time rule, bottom $=0$, apex $=1$, top $=2$, makes the
surface$\leftrightarrow$apex and bottom$\leftrightarrow$top edges the timelike worldlines,
and it stays equivariant.)
\end{remark}

\section{The spin / current sector: Dirac--K\"ahler (\#415)}
\label{sec:dk}
The inhomogeneous cochain $\Omega^\bullet=\bigoplus_k\Omega^k$ is the discrete
Dirac--K\"ahler field, and $D=\dd+\delta$ is its Dirac operator. Assembled from the
integer boundary maps and metric weights (reusing the convention the Hodge Laplacian
already squares), the construction reproduces the expected algebra to machine precision:

\begin{center}
\small
\begin{tabular}{@{}lcl@{}}
\toprule
property & measured & meaning\\
\midrule
$(\dd+\delta)^2 = L$ & residual $\sim 6\times10^{-15}$ & $D$ is a square root of the Hodge Laplacian\\
$\{\gamma^a,\gamma^b\}=2\eta^{ab}$ & residual $0.0$ (exact) & the Clifford algebra of the $1$-cochain action\\
$j^0=\bar\Phi\gamma^0\Phi$ & $=\langle\Phi,\Phi\rangle_W$ & the conserved (Noether) charge density\\
multiplicity & $=4$ & the taste degeneracy $\Lambda(\C^4)=16=4\times4$\\
\bottomrule
\end{tabular}
\end{center}

\noindent The gamma generators are exposed as the $16\times16$ Clifford action of unit
$1$-cochains on the form fiber $\Lambda(\R^4)$ --- the literal realization of ``gammas
$=$ the Clifford action of edges.'' The four-fold multiplicity (the lattice
Dirac--K\"ahler ``doubling'') is the candidate flavor/taste index. The alternative
representation --- the irreducible $4\times4$ Dirac matrices $\gamma^\mu$ --- is recorded
as a documented extension point: it is the image of the same generators under the
Chevalley/K\"ahler--Atiyah isomorphism $\Lambda(\C^4)\to M_4(\C)$, a sibling accessor that
would change only the gamma dimension ($16\to4$) and make the multiplicity an explicit
taste degeneracy; $D$ and $j^0$ are representation-independent. The current's time
component $j^0$ is the charge density; its cross-check against the gauge-field Gauss-law
charge is the subject of \S\ref{sec:charge} and is deferred to \#411.

\section{The charge sector}
\label{sec:charge}

\subsection{Electric and magnetic fields by causal type (\#417)}
The electromagnetic field strength is a $2$-cochain $F=\dd A\in\Omega^2$. On the
Lorentzian complex it splits by the causal type of the plaquette: a plaquette is
\emph{electric} if any of its edges is timelike (a temporal leg), and \emph{magnetic}
if it is purely spacelike. This is implemented as
\texttt{EigenstateSynthesis::curvatureFromConnection} ($F=\dd A$, the same induced-
orientation signed edge sum the periods use) and \texttt{fieldStrengthSplit}, returning
the electric and magnetic parts and their disjoint, complete cell index lists. The tests
confirm the split tracks the causal structure:
\begin{itemize}[nosep]
  \item on a purely spacelike (static) seed: $\norm{E}=0$ exactly, $B=F$ on every cell;
  \item under timelike flux: $\norm{E}>0$, and the electric cells equal the
        independently-recomputed timelike-leg plaquettes;
  \item $E+B=F$ to $10^{-12}$, and $F=\dd A$ is gauge invariant
        ($A\to A+\dd\chi$ leaves $E,B$ unchanged).
\end{itemize}
The reader is purely observational: the proton sector is bit-for-bit unchanged with it
present or absent. Because $E/B$ is read off the cochain by causal type, it is
\emph{interior-shape-independent} --- it holds equally on the prism and the symmetric
apex interior.

\subsection{Electric charge as a Gauss-law holonomy (\#411, in flight)}
The charge ticket reads electric charge as the temporal-sector Gauss law,
$Q=\oint \star F$, a gauged-$U(1)$ holonomy of the field of \S4.1 --- expected to be
metric-robust and quantized, unlike a hand-weighted color covector --- and to verdict the
fractional-($\tfrac13$) charge link to triality. A required cross-check, handed off from
\#415, is that this Gauss-law $Q$ agrees with the Dirac--K\"ahler current's conserved
charge $j^0$ of \S\ref{sec:dk}. This work is in progress; results will be folded in.

\section{The flavor and binding sectors (in flight)}
\label{sec:flight}
Three further pieces are under construction and are reported here only as scope:
\begin{itemize}[nosep]
  \item \textbf{Flavor without a register (\#414, pending).} Seeks the u/d label
  (proton $uud$ $=+1$ vs neutron $udd$ $=0$) in the spatio-temporal split or the
  Dirac--K\"ahler taste multiplicity of \S\ref{sec:dk}, not a parallel hole register. It
  depends on the Gauss-law charge (\#411).
  \item \textbf{The bipartite revisit (\#416, in flight).} Tests whether a twisted
  (orientation-reversing) bipartite tube projects onto the antisymmetric diquark
  $\bar{\rep3}$, and whether a connected two-vertex path reaches the singlet --- on the
  \emph{bipartite/prism} path, which the symmetric apex interior (a tripartite-junction
  construction) does not replace.
  \item \textbf{The $S^3$ dimensional spike (\#418, in flight).} Scopes whether the
  apex/dual-edge octahedron split generalizes to a triangulated $S^3$ base, which
  observables genuinely need $3+1$ D (full $\vec E,\vec B$; the $4\times4$ Dirac matrices)
  versus which are faithful already in $2+1$ D.
\end{itemize}

\section{Faithfulness and no-regression}
Every result above was landed under a fixed set of proton invariants (the track's
$G_1$--$G_8$): the singlet overlap, confinement $\sigmacol\to0$, the exact Stokes charge
conservation, $b_1=11$ with a valid manifold, the color $\Z_3$ spectrum, relabeling-
invariance, determinism, and emergent-first read-out. They held throughout: the geometry
change (\#413) preserved the color sector exactly while sharpening the singlet, the
orientation change (\#412) preserved the manifold and the spectrum while making the
charge a label-independent invariant, and the Dirac--K\"ahler and $E/B$ additions
(\#415, \#417) are observable-only --- the proton observables are unchanged with them
present. No golden result drifted; the twenty junction tests and the per-ticket suites
are green on \texttt{main}.

\section{Assessment and open questions}
\label{sec:assess}
\paragraph{What is established.} The geometry is now honest (no label-sort artifacts),
the proton color singlet is \emph{exact} (overlap $1.000000$, residual $10^{-14}$), the
charge read-out is the geometry's Stokes invariant, the field strength carries a faithful
$E/B$ split, and the Dirac--K\"ahler operator gives genuine $4\times4$-sized Clifford
structure with a conserved current and a four-fold taste index --- all verified to
machine precision and merged.

\paragraph{What is open.} The electric charge as a protected Gauss-law holonomy and its
agreement with $j^0$ (\#411); whether flavor is the taste multiplicity or the
spatio-temporal split (\#414); whether two bipartite steps reach the singlet (\#416);
and whether $3+1$ D over $S^3$ is needed for the full electroweak content (\#418). Two
deeper questions persist from the plan: the cross-sector
\emph{total antisymmetrization} (the $SU(6)$ $\mathbf{56}$ that makes three labels a
proton, not just color $\times$ flavor $\times$ spin), and whether isospin is realized as
a genuine $SU(2)$ with a ladder rather than a static u/d label. Finally,
Remark~\ref{rem:photon} reframes the photon: it should be \emph{sourced}, not
spontaneous, on the symmetric interior --- a concrete next experiment.

\paragraph{Verdict.} The color-confinement skeleton is now exact and the surrounding
machinery (orientation, $E/B$, the Dirac--K\"ahler current) is in place and verified; the
electroweak read-outs that consume it are the work in flight. The track is on its planned
path, with every landed step measured rather than asserted.

\end{document}
